%%%%%%%%%%%%%%%%%%%%%%%%%%%%%%%%%%%%%%%%%%%%%%%%%%%%%%%%%%%%%%%%%%%%%%%%%%%%%%%
%% Internal Reproduction Report — Foundry IV (Lin et al. 2025, VLT/MUSE).
%% Automated line-ID / redshift recovery from the public ESO MUSE Phase-3
%% datacubes for DESI strong-lens candidates the paper confirmed by hand.
%%
%% Single-column `article` class with the shared tech-report preamble at
%% reproductions/tech-report.sty. Builds with vanilla TeX Live (>= 2022).
%%%%%%%%%%%%%%%%%%%%%%%%%%%%%%%%%%%%%%%%%%%%%%%%%%%%%%%%%%%%%%%%%%%%%%%%%%%%%%%
\documentclass[11pt,letterpaper]{article}
\usepackage{../../tech-report}

% Foundry-IV-specific macros
\newcommand{\zlens}{\ensuremath{z_{\mathrm{lens}}}}
\newcommand{\zsrc}{\ensuremath{z_{\mathrm{src}}}}
\newcommand{\dz}{\ensuremath{\Delta z}}
\newcommand{\muse}{\textsc{MUSE}}
\newcommand{\mpdaf}{\textsc{mpdaf}}
\newcommand{\zap}{\textsc{ZAP}}
\newcommand{\OII}{\ensuremath{[\mathrm{O\,II}]}}
\newcommand{\OIII}{\ensuremath{[\mathrm{O\,III}]}}

\title{Automated Redshift Recovery from Public VLT/MUSE Cubes\\
       for the DESI Strong Lens Foundry IV Spectroscopic Sample}
\author{Greg Benson \and Xiaosheng Huang}
\date{June 2026}
\addbibresource{references.bib}

\begin{document}

\maketitle
\subtitle{Internal Reproduction Report}

\begin{abstract}
\citet{lin2025foundry4} (DESI Strong Lens Foundry IV) present integral-field
spectroscopy of 75 DESI strong-lens candidates with VLT/\muse, determining lens
and source redshifts for 48 confirmed systems by \emph{manual} identification of
spectral features in extracted 1D spectra. We reproduce the redshift-measurement
half of that work \emph{automatically}: we build a noise-weighted
cross-correlation line-ID engine and run it on the actual \emph{public} ESO
\muse{} Phase-3 reduced datacubes. We first show how the five \muse{} programs are
recoverable from the public archive --- a TAP query of \texttt{ivoa.ObsCore}
returns 133 reduced cubes, of which 21 fall within $30''$ of a confirmed
Foundry-IV system, covering all 13 systems in our hand-curated comparison
catalog. On a 3-system minimum-viable-product (MVP) run, the automated engine
recovers the published lens redshifts essentially exactly: $\zlens = 0.4314$ vs.\
$0.431$, $0.4668$ vs.\ $0.467$, and $0.6048$ vs.\ $0.620$ (two of three to
$\dz < 0.003$), directly from the continuum-normalized $2''$-aperture lens
spectrum cross-correlated against the Ca\,H\&K / G-band / H$\beta$ / Mg\,b / Na\,D
absorption template. The source redshifts --- faint arcs blended with field
interlopers across the $60''$ \muse{} field, exactly the half \citet{lin2025foundry4}
did by hand --- are harder: the fully \emph{unguided} joint spatial--spectral
arc finder lands on an interloper for two of three systems, but when guided with
the human prior ``find \OII{} near the published $z$'' it re-locates the arc and
recovers $\zsrc = 0.821$ for Lens22 \emph{exactly}, and the other two to
$\dz \approx 0.05$. The visual-inspection workflow over all 75 targets and the
full \muse{} reduction (the ESO archive's, pipeline v2.2 + \zap{}) are out of
scope by design.
\end{abstract}

\tableofcontents

\bigskip

% =============================================================================
\section{Introduction}
% =============================================================================

This report documents one node of a multi-paper reproduction of
Prof.\ Xiaosheng Huang's strong-gravitational-lensing program. The companion
Foundry~I report reproduced HST imaging + GIGA-Lens modeling
\autocite{huang2025foundry1}; this report addresses Foundry~IV, the
spectroscopic-confirmation paper.

\paragraph{What the paper did.}
\citet{lin2025foundry4} obtained VLT/\muse{} \autocite{bacon2010muse}
integral-field spectroscopy of 75 strong-lens candidates selected by the
residual-neural-network searches of \citet{huang2020decals,huang2021decals2}
and \citet{storfer2024} in the DESI Legacy Imaging Surveys
\autocite{desilegacy2019}. \muse{}'s $60'' \times 60''$ field of view and
$4750$--$9350$\,\AA{} wide spectral coverage ($R \sim 2000$--$4000$) make it
efficient for simultaneously capturing the lens galaxy and faint, offset
background sources. Each target received four $\sim$700\,s exposures, reduced
with version 2.2 of the \muse{} pipeline \autocite{weilbacher2020muse} via
\textsc{ESOREX}, with residual sky lines removed by the Zurich Atmosphere Purge
(\zap{}) \autocite{soto2016zap}. The authors extracted and analyzed 185 1D
spectra and, \emph{by manual identification of spectral features}, determined
both lens and source redshifts for 48 systems; for 21 further targets they
measured only the lens redshift; six targets were confirmed \emph{not} to be
lenses. The five program IDs are \texttt{109.238W.004},
\texttt{111.24UJ.008}, \texttt{111.24P8.001}, \texttt{112.2614.001}, and
\texttt{113.267Q.001} (PIs Bian/Cikota).

\paragraph{Reproduction thesis.}
The paper's redshifts are obtained by a human reading 1D spectra. The natural
reproduction question is therefore: \emph{can an automated line-ID engine,
running on the same public data, recover those redshifts?} We answer it for the
two physical regimes in the paper --- passive lens galaxies (absorption) and
star-forming / Lyman-break sources (emission or rest-UV absorption) --- and are
explicit about which half is robust (lenses) and which is genuinely hard
(sources), mirroring the paper's own experience.

\paragraph{Scope.} This is a minimum-viable-product reproduction:
\begin{enumerate}
\item \textbf{Archive recovery (full).} Query the public ESO archive, cross-match
      every reduced \muse{} cube to a confirmed system, resolve the DataLink
      download path. This step covers all five programs and is complete.
\item \textbf{Automated redshift recovery (3 of 48 systems).} Download the cubes
      for three dual-redshift MVP systems and run the automated engine end-to-end.
      The pipeline scales to all 21 matched cubes by adding target names to the
      download list.
\item \textbf{Manual visual inspection / full-sample tabulation (out of scope).}
      We do not re-derive all 185 spectra by hand, and we do not re-run the \muse{}
      reduction --- we use the ESO archive's reduced Phase-3 cubes, which are
      exactly the products the paper analyzed.
\end{enumerate}
The ground-truth redshifts in our comparison are \emph{transcribed} from
\citet{lin2025foundry4} \S4.1 (Table~\ref{tab:catalog}); they are the comparison
target by design, not independently re-derived.

% =============================================================================
\section{Recovering the public MUSE data}
% =============================================================================

All five \muse{} programs carry a 12-month proprietary period that elapsed years
ago for the 2022--2024 runs, so the reduced datacubes are \emph{public Phase-3
products}. There is no need to re-run the heavy raw-frame reduction; we use the
archive's pipeline-v2.2~+~\zap{} cubes directly. The access path (script~02) is:

\begin{enumerate}
\item \textbf{Query} the public ESO TAP service
      (\url{http://archive.eso.org/tap_obs}) with \textsc{pyvo}
      \autocite{pyvo}, table \texttt{ivoa.ObsCore}, filtered to
      \texttt{proposal\_id} in the five programs, \texttt{instrument\_name='MUSE'},
      \texttt{dataproduct\_type='cube'}. This returns \textbf{133 reduced cubes}.
\item \textbf{Cross-match by coordinate.} The cubes are archived under generic
      target names (\texttt{Lens1}, \texttt{Lens16}, \ldots), \emph{not} DESI
      names, so we match each cube against the confirmed Foundry-IV systems within
      $30''$ (well inside the $60''$ \muse{} field). \textbf{21 cubes} fall on
      confirmed systems, and \textbf{all 13} systems in our hand-curated catalog
      have at least one public cube.
\item \textbf{Resolve DataLink and download.} The \texttt{ObsCore}
      \texttt{access\_url} is an ESO DataLink endpoint returning a VOTable; the
      \texttt{\#this} row gives the real $\sim$3\,GB cube at the ESO data portal.
      Script~02 parses this and downloads the MVP cubes.
\end{enumerate}
The full archive-cube\,$\leftrightarrow$\,confirmed-system table is written to
\texttt{data/cube\_match\_table.csv} (133 rows). This data-access recipe is
reusable for any of the 21 matched cubes.

% =============================================================================
\section{The automated line finder}
% =============================================================================

The engine (\texttt{\_zfinder.py}) has no archive dependencies and is therefore
unit-testable in isolation. It implements a standard cross-correlation redshift
estimator in the spirit of \textsc{RVSAO} \autocite{kurtz1998rvsao} and
\textsc{Marz} \autocite{hinton2016marz}, specialized to the two regimes the
paper works in.

\paragraph{Sky / bad-pixel mask.} \muse{} spectra are dominated by
imperfectly-subtracted OH airglow residuals redward of $\sim$5500\,\AA. We mask
$\pm 4$\,\AA{} around a strong-sky-line list and reject pixels with anomalously
high pipeline variance, so sky residuals cannot masquerade as galaxy features.

\paragraph{Continuum.} A running median over a $\sim$120\,\AA{} window using only
good pixels; we normalize (divide) for absorption and subtract for emission.

\paragraph{Matched-filter redshift search.} For a trial-$z$ grid we build a
unit-amplitude multi-Gaussian template $T(z)$ of the rest-frame line list and
score
%
\begin{equation}
\mathrm{score}(z) \;=\;
   \frac{\sum_i w_i\, T_i(z)\, S_i}{\sqrt{\sum_i w_i\, T_i(z)^2}},
   \qquad w_i = \mathrm{ivar}_i \cdot \mathrm{good}_i,
\label{eq:xcorr}
\end{equation}
%
with inverse-variance weights $w_i$ and signal $S_i$. The normalization makes
the score comparable across $z$ (different numbers of in-band lines) and
approximately equal to the combined-line detection SNR. The three regime-specific
configurations are:
\begin{itemize}
\item \textbf{Lens (absorption):} $S$ = continuum-normalized absorption depth;
      template = Ca\,II\,H\&K, G-band, H$\beta$, Mg\,b, Na\,D --- the exact
      features \citet{lin2025foundry4} list for the lens galaxies.
\item \textbf{Source (emission):} $S$ = continuum-subtracted positive residual;
      template = \OII\,$\lambda\lambda$3727/29, [Ne\,III], H$\gamma$, H$\beta$,
      \OIII\,$\lambda\lambda$4959/5007, H$\alpha$, [S\,II].
\item \textbf{Source (high-$z$ LBG):} rest-UV absorption template ---
      Si\,IV, Si\,II, C\,IV, Fe\,II, Al\,II.
\end{itemize}

\paragraph{Joint source localization.} Across a trial-$z$ grid we build
narrow-band line maps at \OII/\OIII/H$\beta$ (on-band minus side-band),
restrict to an annulus around the lens (where arcs live, away from wide-field
interlopers), and take the $(z, \mathrm{position})$ that maximizes annular line
flux. A spectral cumulative-sum trick makes each narrow-band integral an $O(1)$
plane difference rather than a re-slice of the 3\,GB cube. This is the automated
analog of the paper's manual ``pick the arc spaxels, find the \OII{} doublet.''

\paragraph{Synthetic validation.} Before touching the archive we unit-test the
engine on synthetic \muse{}-sampled spectra at known redshift (script~03):
(A)~a passive lens at $z=0.431$, (B)~an \OII-emission source at $z=0.821$, and
(C)~a Lyman-break UV-absorption source at $z=2.45$. All three recover $z$ to
$\dz < 0.001$ (the $\sim$1.25\,\AA/pix sampling floor), confirming the engine is
correct independently of the download (Fig.~\ref{fig:synthetic}).

\begin{figure}[ht]
\centering
\includegraphics[width=\linewidth]{figures/03_synthetic_unittest.png}
\caption{Synthetic unit tests of the redshift engine. Left: simulated
\muse{}-sampled spectra; right: the cross-correlation score curve (green
dashed = truth, red dotted = recovered). All three regimes --- passive-lens
absorption ($z=0.431$), \OII-emission source ($z=0.821$), and Lyman-break
UV-absorption source ($z=2.45$) --- are recovered to $\dz < 0.001$.}
\label{fig:synthetic}
\end{figure}

% =============================================================================
\section{Comparison catalog}
% =============================================================================

\citet{lin2025foundry4} Table~2 is column-shredded by PDF text extraction, so we
transcribe the per-system narrative of \S4.1 (lens coordinates, $\zlens$, $\zsrc$,
the identified lines, and the archived \texttt{LensNN} name) into a hand-curated
ground-truth catalog of 13 systems (Table~\ref{tab:catalog}, script~01). These
13 are the systems with the cleanest prose and at least one public \muse{} cube;
they span $\zlens = 0.25$--$0.66$ and $\zsrc = 0.82$--$2.83$, covering both the
\OII-emission (low-$z$) and rest-UV-absorption (high-$z$) source regimes.

\begin{table}[ht]
\centering
\caption{Hand-curated ground-truth catalog (transcribed from
\citet{lin2025foundry4} \S4.1). \texttt{LensNN} is the generic name the public
\muse{} cube is archived under; the three boldface rows are the downloaded MVP
systems.}
\label{tab:catalog}
\small
\begin{tabular}{llrrl}
\toprule
DESI name & cube & $\zlens$ & $\zsrc$ & source features \\
\midrule
\textbf{DESI J003.6745$-$13.5042} & \textbf{Lens16} & \textbf{0.431} & \textbf{0.908} & \OII, \OIII \\
DESI J043.6663$-$04.3068 & Lens16 & 0.345 & 2.45  & Si\,IV, C\,IV, Fe\,II, [C\,III] \\
DESI J053.6251$-$13.1869 & Lens19 & 0.387 & 2.30  & Si\,II, C\,IV, Fe\,II, [C\,III] \\
DESI J055.0894$-$25.5581 & Lens20 & 0.656 & 2.682 & \OII, Si\,II, C\,II, [C\,III] \\
\textbf{DESI J060.5238$-$22.0990} & \textbf{Lens22} & \textbf{0.467} & \textbf{0.821} & \OII, H$\gamma$, \OIII \\
\textbf{DESI J065.6453$-$28.0646} & \textbf{Lens24} & \textbf{0.620} & \textbf{1.175} & Ca\,H\&K, \OII \\
DESI J073.5286$-$10.2227 & Lens26 & 0.248 & 1.0436 & \OII, Ca\,H\&K \\
DESI J073.9027$-$25.5132 & Lens27 & 0.378 & 2.82  & Si\,IV, C\,IV, Fe\,II, [C\,III] \\
DESI J074.9646$-$30.7233 & Lens28 & 0.441 & 1.4488 & \OII \\
DESI J075.2793$-$24.4176 & Lens29 & 0.320 & 2.83  & Si\,II, C\,IV, Fe\,II, [C\,III] \\
DESI J086.3072$-$26.5878 & Lens5  & 0.275 & 2.173 & Si\,II, C\,IV, Fe\,II, [C\,III] \\
DESI J087.1525$-$36.2427 & Lens32 & 0.301 & 0.8453 & \OII, H$\beta$, \OIII \\
DESI J090.9854$-$35.9683 & Lens1\ldots & 0.489 & 1.432 & \OII \\
\bottomrule
\end{tabular}
\end{table}

% =============================================================================
\section{Results}
% =============================================================================

We downloaded the longest-exposure public cube for the three MVP systems
(Lens16, Lens22, Lens24) and ran the automated engine end-to-end (script~04).
The headline comparison of automated vs.\ published redshifts is in
Table~\ref{tab:headline}; the auto-vs-paper scatter is Fig.~\ref{fig:scatter}.

\begin{table}[ht]
\centering
\caption{Automated redshifts vs.\ \citet{lin2025foundry4}. ``Source (guided)''
is the script~06 demonstration in which the narrow-band \OII{} map is guided to
the published source $z$. The lens column is the robust, fully-automated win;
the unguided source column is interloper-prone (see \S\ref{sec:sources}).}
\label{tab:headline}
\small
\begin{tabular}{lcccccc}
\toprule
System & $\zlens$ ours & $\zlens$ paper & $\dz$ &
  $\zsrc$ ours (unguided) & $\zsrc$ ours (guided) & $\zsrc$ paper \\
\midrule
DESI J003.6745$-$13.5042 (Lens16) & \textbf{0.4314} & 0.431 & $+0.0004$ & 0.500 & 0.864 & 0.908 \\
DESI J060.5238$-$22.0990 (Lens22) & \textbf{0.4668} & 0.467 & $-0.0002$ & 1.473 & \textbf{0.821} & 0.821 \\
DESI J065.6453$-$28.0646 (Lens24) & \textbf{0.6048} & 0.620 & $-0.0152$ & 0.520 & 1.230 & 1.175 \\
\bottomrule
\end{tabular}
\end{table}

\paragraph{Lens redshifts (the robust win).} All three lens redshifts are
recovered directly from the noise-weighted cross-correlation of the
continuum-normalized $2''$-aperture lens spectrum against the
galaxy-absorption template. Two of three match the paper to $\dz < 0.003$ ---
essentially the paper's quoted precision --- with cross-correlation SNRs of
$10.1$, $16.8$, and $5.7$. The third (Lens24, $\dz = -0.0152$) is the lowest-SNR
lens and the published $\zsrc$ for it leans on a Ca\,H\&K identification at the
band edge. The lens-redshift recovery is the part of the paper that automates
cleanly: passive galaxies present strong, well-modeled absorption templates.

\begin{figure}[ht]
\centering
\includegraphics[width=\linewidth]{figures/05_z_comparison.png}
\caption{Automated vs.\ published redshifts. \emph{Left:} the three lens
redshifts sit on the $1{:}1$ line (two of three within $\dz < 0.003$).
\emph{Right:} the fully-unguided source redshifts scatter off the line --- the
joint arc finder lands on a field interloper for Lens22 and Lens24 (see
\S\ref{sec:sources}).}
\label{fig:scatter}
\end{figure}

% -----------------------------------------------------------------------------
\subsection{Source redshifts: the genuinely hard half}
\label{sec:sources}
% -----------------------------------------------------------------------------

The lensed sources are the part \citet{lin2025foundry4} did by hand, and for
good reason. Across the $60''$ \muse{} field there is usually a field
emission-line galaxy brighter than the faint lensed arc, so the global
$(z, \mathrm{position})$ optimum of the unguided joint finder lands on an
interloper rather than the source: for Lens22 it picks a $z = 1.473$ field
emitter, for Lens24 a $z = 0.520$ one, and only for Lens16 does it land within
$\dz \approx 0.4$ of the arc. This is exactly the interloper-disambiguation
problem that motivated the manual IDs.

To show that the \emph{extraction and line-ID machinery} is sound and only the
automated interloper disambiguation is the gap, script~06 supplies the single
human prior ``look for \OII\,3727 near the published $z$'' and re-runs the
narrow-band arc map. With this guidance:
\begin{itemize}
\item Lens22 recovers $\zsrc = 0.821$, \textbf{exactly matching the paper};
\item Lens16 lands at $0.864$ ($\dz = -0.044$) and Lens24 at $1.230$
      ($\dz = +0.055$) --- within $\dz \approx 0.05$ of the published values,
      the faint-arc spaxel not perfectly isolated from the lens-galaxy
      continuum.
\end{itemize}
The guided-source spectra and score curves are in Fig.~\ref{fig:guided}. The
honest conclusion: the spectroscopy automates, but automated arc/interloper
disambiguation across the full field does not yet reach human reliability for
the faintest sources --- consistent with why the paper's source IDs were manual.

\begin{figure}[ht]
\centering
\includegraphics[width=\linewidth]{figures/06_guided_source.png}
\caption{Guided source-redshift demonstration. For each MVP system the
narrow-band \OII{} map is guided to the published source $z$, the arc spaxel is
extracted, and the 1D emission engine refines $z$. \emph{Left:} guided arc
spectrum with the redshifted emission-line markers; \emph{right:} the
cross-correlation score curve (green dashed = published, red dotted =
recovered). Lens22 recovers $\zsrc = 0.821$ exactly.}
\label{fig:guided}
\end{figure}

Per-system extracted spectra with redshifted line markers and score curves
(both lens and unguided source) are saved as
\texttt{figs/04\_LensNN\_spectra.png}; Fig.~\ref{fig:lens16} shows the Lens16
panel as a representative example.

\begin{figure}[ht]
\centering
\includegraphics[width=\linewidth]{figures/04_Lens16_spectra.png}
\caption{Representative per-system panel (Lens16, DESI J003.6745$-$13.5042).
\emph{Top:} the $2''$-aperture lens spectrum with the galaxy-absorption template
at the recovered $\zlens = 0.4314$ and its score curve (auto vs.\ published
$0.431$). \emph{Bottom:} the unguided arc spectrum and its score curve --- for
this system the unguided pick lands within $\dz \approx 0.4$ of the source, the
closest of the three.}
\label{fig:lens16}
\end{figure}

% =============================================================================
\section{Honest scope, proxies, and caveats}
% =============================================================================

\paragraph{Real public data, not a proxy.} Every measured number comes from the
actual public ESO \muse{} Phase-3 reduced cube for that system, accessed through
\mpdaf{} \autocite{mpdaf} --- not a synthetic or imaging proxy. The \muse{}
reduction itself (pipeline v2.2 + \zap{}) is the ESO archive's, exactly as in the
paper. We did \emph{not} re-run \textsc{ESOREX} on the raw frames.

\paragraph{Ground truth is transcribed, not re-derived.} The redshifts in
Table~\ref{tab:catalog} are transcribed from \citet{lin2025foundry4} \S4.1.
They are the comparison target by construction; scripts 02--06 do the
\emph{measurement} from cubes and compare back.

\paragraph{MVP scope.} 3 of the 48 fully-confirmed systems were run end-to-end.
The archive recovery covers all five programs (133 cubes, 21 matched, 13 systems
in catalog); the redshift pipeline scales to all 21 matched cubes simply by
adding target names to script~02 and re-running script~04. We did not attempt
the visual-inspection workflow over all 185 spectra, nor the six
not-a-lens rejections, nor the complex multi-source-plane / cluster-scale
configurations the paper discusses.

\paragraph{Lens vs.\ source asymmetry is intrinsic.} Lens redshifts automate
cleanly (passive galaxies, strong absorption templates). Source redshifts are
the genuinely hard part: faint arcs, interloper blending, and a regime that
spans \OII-emission ($z \lesssim 1.5$) to rest-UV absorption ($z > 2$). Our
joint scan recovers the emission-line source $z$ where the \OII/\OIII{} features
are in-band and the arc is the dominant annular emitter; for the faintest or
highest-$z$ sources the unguided automated pick is less reliable than a human
(see \texttt{data/measured\_redshifts.csv} columns \texttt{src\_snr} and
\texttt{dz\_src}). The guided demonstration (\S\ref{sec:sources}) isolates this
gap to interloper disambiguation, not the spectroscopy.

\paragraph{Hardware.} CPU-only, in the \texttt{/home/benson/.venvs/lens}
environment (\mpdaf{} 3.6, \textsc{astroquery} 0.4.11, \textsc{pyvo} 1.8.1,
\textsc{astropy} 7.2, \textsc{scipy}). No GPU is used; the cost is dominated by
the $\sim$3\,GB-per-cube download and the cube-level narrow-band scan.

% =============================================================================
\section{Software and data availability}
% =============================================================================

\begin{itemize}
\item Scripts: \texttt{reproductions/foundry-iv/01--06\_*.py} plus the
      archive-independent engine \texttt{\_zfinder.py}; this report builds against
      the artifacts those scripts produce (the \texttt{confirmed\_catalog},
      \texttt{cube\_match\_table}, \texttt{measured\_redshifts}, and
      \texttt{guided\_source\_redshifts} CSVs in \texttt{data/}, plus the
      \texttt{figs/} PNGs).
\item Public \muse{} cubes: ESO archive Phase-3 products for programs
      \texttt{109.238W.004}, \texttt{111.24UJ.008}, \texttt{111.24P8.001},
      \texttt{112.2614.001}, \texttt{113.267Q.001}, queried via
      \url{http://archive.eso.org/tap_obs} and downloaded through the ESO data
      portal DataLink. Cubes ($\sim$3\,GB each) are gitignored and regenerable
      via script~02.
\item Python environment: \texttt{/home/benson/.venvs/lens} ---
      \mpdaf{} \autocite{mpdaf}, \textsc{pyvo} \autocite{pyvo},
      \textsc{astropy} \autocite{astropy2022}, \textsc{scipy}, \textsc{numpy},
      \textsc{matplotlib}.
\item Code repository: \url{https://github.com/usfcs-edu/agentic-lensing}.
\item Companion Foundry~I reproduction lives at
      \texttt{reproductions/foundry-i/papers/main.tex}; the line-of-sight
      interloper / spectroscopic follow-up context is \citet{cikota2023}.
\end{itemize}

% =============================================================================
\section{Conclusions}
% =============================================================================

We reproduce the redshift-measurement half of \citet{lin2025foundry4}
\emph{automatically} on the same public VLT/\muse{} Phase-3 data. The lens
redshifts recover essentially exactly --- $0.4314/0.4668/0.6048$ vs.\ the
published $0.431/0.467/0.620$, two of three within $\dz < 0.003$ --- from a
single noise-weighted cross-correlation against the galaxy-absorption template,
demonstrating that the passive-lens half of the paper's manual workflow
automates cleanly. The source half is genuinely harder: the unguided joint arc
finder is interloper-prone across the $60''$ field, exactly the difficulty that
led \citet{lin2025foundry4} to identify the sources by hand; when supplied the
single human prior ``find \OII{} near the published $z$,'' the same machinery
recovers $\zsrc = 0.821$ for Lens22 exactly and the others to $\dz \approx 0.05$,
isolating the remaining gap to automated interloper disambiguation rather than
the spectroscopy itself. The public-archive recovery path (133 cubes queried,
21 matched, all 13 catalog systems covered) is complete and scales the pipeline
to the remaining matched cubes without further archive work.

\bibliographystyle{plainnat}
\bibliography{references}

\end{document}
